%%%%%%%%%%%%%%%%%%%%%%%%%%%%%%%%%%%%%%%%%%%%%%%%%%%%%%%%%%%%%%%%%%%%%%%%%%%%%%%%
% CHAPTER 7
%%%%%%%%%%%%%%%%%%%%%%%%%%%%%%%%%%%%%%%%%%%%%%%%%%%%%%%%%%%%%%%%%%%%%%%%%%%%%%%%

\chapter{Discrete-Time Signals and Systems}

\begin{tcolorbox}[colback=blue!5!white,colframe=blue!70!black,title=Learning Objectives]

After studying this chapter, the reader will be able to

\begin{itemize}
\item Define discrete-time signals and sequences.
\item Understand basic sequence operations.
\item Define discrete-time systems.
\item Test linearity, time invariance, causality and stability.
\item Understand discrete-time convolution.
\item Compute the output of an LTI discrete-time system.
\end{itemize}

\end{tcolorbox}

%%%%%%%%%%%%%%%%%%%%%%%%%%%%%%%%%%%%%%%%%%%%%%%%%%%%%%%%%%%%%%%%%%%%%%%%%%%%%%%%

\section{Introduction}

A \textbf{discrete-time signal} is a signal defined only at discrete instants of time.

Instead of \(x(t)\), we use the notation

\[
\boxed{
x[n]
}
\]

where \(n\) is an integer.

Discrete-time signals arise naturally from sampling continuous-time signals and are
the basic signals processed by digital computers, DSP processors and microcontrollers.

%%%%%%%%%%%%%%%%%%%%%%%%%%%%%%%%%%%%%%%%%%%%%%%%%%%%%%%%%%%%%%%%%%%%%%%%%%%%%%%%

\section{Continuous-Time vs Discrete-Time}

\begin{table}[H]
\centering
\caption{Comparison of signal types}
\begin{tabular}{|p{4cm}|p{4cm}|p{4cm}|}
\hline
\textbf{Feature} & \textbf{Continuous-Time} & \textbf{Discrete-Time}\\
\hline
Independent variable & \(t\) & \(n\)\\
\hline
Definition & All real values & Integer values only\\
\hline
Notation & \(x(t)\) & \(x[n]\)\\
\hline
Example & \(\sin(2\pi t)\) & \(\sin(0.2\pi n)\)\\
\hline
\end{tabular}
\end{table}

%%%%%%%%%%%%%%%%%%%%%%%%%%%%%%%%%%%%%%%%%%%%%%%%%%%%%%%%%%%%%%%%%%%%%%%%%%%%%%%%

\section{Sampling Relationship}

If a continuous-time signal is sampled every \(T_s\) seconds, then

\[
\boxed{
x[n]=x(nT_s)
}
\]

where

\begin{itemize}
\item \(T_s\): sampling period,
\item \(f_s=1/T_s\): sampling frequency.
\end{itemize}

%%%%%%%%%%%%%%%%%%%%%%%%%%%%%%%%%%%%%%%%%%%%%%%%%%%%%%%%%%%%%%%%%%%%%%%%%%%%%%%%

\section{Basic Sequences}

%%%%%%%%%%%%%%%%%%%%%%%%%%%%%%%%%%%%%%%%%%%%%%%%%%%%%%%%%%%%%%%%%%%%%%%%%%%%%%%%

\subsection{Unit Sample Sequence}

\[
\boxed{
\delta[n]=
\begin{cases}
1, & n=0\\
0, & n\neq0
\end{cases}
}
\]

%%%%%%%%%%%%%%%%%%%%%%%%%%%%%%%%%%%%%%%%%%%%%%%%%%%%%%%%%%%%%%%%%%%%%%%%%%%%%%%%

\subsection{Unit Step Sequence}

\[
\boxed{
u[n]=
\begin{cases}
1, & n\ge0\\
0, & n<0
\end{cases}
}
\]

%%%%%%%%%%%%%%%%%%%%%%%%%%%%%%%%%%%%%%%%%%%%%%%%%%%%%%%%%%%%%%%%%%%%%%%%%%%%%%%%

\subsection{Exponential Sequence}

\[
\boxed{
x[n]=a^n u[n]
}
\]

\begin{itemize}
\item \(|a|<1\): decaying,
\item \(|a|>1\): growing,
\item \(a<0\): alternating sign.
\end{itemize}

%%%%%%%%%%%%%%%%%%%%%%%%%%%%%%%%%%%%%%%%%%%%%%%%%%%%%%%%%%%%%%%%%%%%%%%%%%%%%%%%

\section{Sequence Operations}

%%%%%%%%%%%%%%%%%%%%%%%%%%%%%%%%%%%%%%%%%%%%%%%%%%%%%%%%%%%%%%%%%%%%%%%%%%%%%%%%

\subsection{Time Shifting}

Delay by \(n_0\):

\[
\boxed{
x[n-n_0]
}
\]

Advance by \(n_0\):

\[
\boxed{
x[n+n_0]
}
\]

%%%%%%%%%%%%%%%%%%%%%%%%%%%%%%%%%%%%%%%%%%%%%%%%%%%%%%%%%%%%%%%%%%%%%%%%%%%%%%%%

\subsection{Time Reversal}

\[
\boxed{
x[-n]
}
\]

%%%%%%%%%%%%%%%%%%%%%%%%%%%%%%%%%%%%%%%%%%%%%%%%%%%%%%%%%%%%%%%%%%%%%%%%%%%%%%%%

\subsection{Amplitude Scaling}

\[
\boxed{
y[n]=A\,x[n]
}
\]

%%%%%%%%%%%%%%%%%%%%%%%%%%%%%%%%%%%%%%%%%%%%%%%%%%%%%%%%%%%%%%%%%%%%%%%%%%%%%%%%

\section{Even and Odd Sequences}

%%%%%%%%%%%%%%%%%%%%%%%%%%%%%%%%%%%%%%%%%%%%%%%%%%%%%%%%%%%%%%%%%%%%%%%%%%%%%%%%

\subsection{Even Sequence}

\[
\boxed{
x[n]=x[-n]
}
\]

%%%%%%%%%%%%%%%%%%%%%%%%%%%%%%%%%%%%%%%%%%%%%%%%%%%%%%%%%%%%%%%%%%%%%%%%%%%%%%%%

\subsection{Odd Sequence}

\[
\boxed{
x[n]=-x[-n]
}
\]

%%%%%%%%%%%%%%%%%%%%%%%%%%%%%%%%%%%%%%%%%%%%%%%%%%%%%%%%%%%%%%%%%%%%%%%%%%%%%%%%

\subsection{Decomposition}

Any sequence can be written as

\[
\boxed{
x[n]=x_e[n]+x_o[n]
}
\]

where

\[
\boxed{
x_e[n]=\frac{x[n]+x[-n]}{2}
}
\]

and

\[
\boxed{
x_o[n]=\frac{x[n]-x[-n]}{2}
}
\]

%%%%%%%%%%%%%%%%%%%%%%%%%%%%%%%%%%%%%%%%%%%%%%%%%%%%%%%%%%%%%%%%%%%%%%%%%%%%%%%%

\section{Periodic Discrete-Time Signals}

A sequence is periodic if

\[
\boxed{
x[n+N]=x[n]
}
\]

for some positive integer \(N\).

The smallest such integer is the \textbf{fundamental period}.

%%%%%%%%%%%%%%%%%%%%%%%%%%%%%%%%%%%%%%%%%%%%%%%%%%%%%%%%%%%%%%%%%%%%%%%%%%%%%%%%

\subsection{Sinusoidal Sequence}

Consider

\[
x[n]=\cos(\omega_0 n).
\]

It is periodic if

\[
\boxed{
\frac{\omega_0}{2\pi}=\frac{m}{N}
}
\]

where \(m\) and \(N\) are integers.

%%%%%%%%%%%%%%%%%%%%%%%%%%%%%%%%%%%%%%%%%%%%%%%%%%%%%%%%%%%%%%%%%%%%%%%%%%%%%%%%

\subsection{Example}

\[
x[n]=\cos\left(\frac{\pi}{4}n\right)
\]

Here

\[
\frac{\omega_0}{2\pi}=\frac{1}{8}.
\]

Hence,

\[
\boxed{
N=8
}
\]

%%%%%%%%%%%%%%%%%%%%%%%%%%%%%%%%%%%%%%%%%%%%%%%%%%%%%%%%%%%%%%%%%%%%%%%%%%%%%%%%

\section{Discrete-Time Systems}

A discrete-time system maps an input sequence \(x[n]\) to an output sequence
\(y[n]\).

\[
\boxed{
y[n]=T\{x[n]\}
}
\]

%%%%%%%%%%%%%%%%%%%%%%%%%%%%%%%%%%%%%%%%%%%%%%%%%%%%%%%%%%%%%%%%%%%%%%%%%%%%%%%%

\section{Linearity}

A system is linear if

\[
T\{a_1x_1[n]+a_2x_2[n]\}
=
a_1T\{x_1[n]\}+a_2T\{x_2[n]\}.
\]

%%%%%%%%%%%%%%%%%%%%%%%%%%%%%%%%%%%%%%%%%%%%%%%%%%%%%%%%%%%%%%%%%%%%%%%%%%%%%%%%

\subsection{Example}

\[
y[n]=3x[n]-2x[n-1]
\]

is linear.

%%%%%%%%%%%%%%%%%%%%%%%%%%%%%%%%%%%%%%%%%%%%%%%%%%%%%%%%%%%%%%%%%%%%%%%%%%%%%%%%

\section{Time Invariance}

A system is time invariant if delaying the input delays the output by the same amount.

If

\[
x[n]\rightarrow y[n],
\]

then

\[
x[n-n_0]\rightarrow y[n-n_0].
\]

%%%%%%%%%%%%%%%%%%%%%%%%%%%%%%%%%%%%%%%%%%%%%%%%%%%%%%%%%%%%%%%%%%%%%%%%%%%%%%%%

\subsection{Example}

\[
y[n]=x[n]+x[n-1]
\]

is time invariant.

%%%%%%%%%%%%%%%%%%%%%%%%%%%%%%%%%%%%%%%%%%%%%%%%%%%%%%%%%%%%%%%%%%%%%%%%%%%%%%%%

\section{Causality}

A system is causal if the output at time \(n\) depends only on present and past inputs.

%%%%%%%%%%%%%%%%%%%%%%%%%%%%%%%%%%%%%%%%%%%%%%%%%%%%%%%%%%%%%%%%%%%%%%%%%%%%%%%%

\subsection{Causal}

\[
y[n]=x[n]+x[n-2]
\]

%%%%%%%%%%%%%%%%%%%%%%%%%%%%%%%%%%%%%%%%%%%%%%%%%%%%%%%%%%%%%%%%%%%%%%%%%%%%%%%%

\subsection{Noncausal}

\[
y[n]=x[n+1]
\]

because it depends on a future input.

%%%%%%%%%%%%%%%%%%%%%%%%%%%%%%%%%%%%%%%%%%%%%%%%%%%%%%%%%%%%%%%%%%%%%%%%%%%%%%%%

\section{BIBO Stability}

A discrete-time system is stable if every bounded input produces a bounded output.

For an LTI system with impulse response \(h[n]\),

\[
\boxed{
\sum_{n=-\infty}^{\infty}|h[n]|<\infty
}
\]

%%%%%%%%%%%%%%%%%%%%%%%%%%%%%%%%%%%%%%%%%%%%%%%%%%%%%%%%%%%%%%%%%%%%%%%%%%%%%%%%

\section{Impulse Response}

The response to the unit sample sequence is called the impulse response.

\[
\boxed{
h[n]=T\{\delta[n]\}
}
\]

For LTI systems, \(h[n]\) completely characterizes the system.

%%%%%%%%%%%%%%%%%%%%%%%%%%%%%%%%%%%%%%%%%%%%%%%%%%%%%%%%%%%%%%%%%%%%%%%%%%%%%%%%

\section{Convolution Sum}

For a discrete-time LTI system,

\[
\boxed{
y[n]=x[n]*h[n]
}
\]

where

\[
\boxed{
y[n]=
\sum_{k=-\infty}^{\infty}
x[k]\,h[n-k]
}
\]

This is the discrete-time convolution sum.

%%%%%%%%%%%%%%%%%%%%%%%%%%%%%%%%%%%%%%%%%%%%%%%%%%%%%%%%%%%%%%%%%%%%%%%%%%%%%%%%

\section{Graphical Steps for Convolution}

To compute \(y[n]\):

\begin{enumerate}
\item Flip \(h[k]\rightarrow h[-k]\),
\item Shift \(h[-k]\rightarrow h[n-k]\),
\item Multiply \(x[k]\) and \(h[n-k]\),
\item Sum over all \(k\).
\end{enumerate}

%%%%%%%%%%%%%%%%%%%%%%%%%%%%%%%%%%%%%%%%%%%%%%%%%%%%%%%%%%%%%%%%%%%%%%%%%%%%%%%%

\section{Example 1}

Let

\[
x[n]=u[n],
\qquad
h[n]=u[n].
\]

Then

\[
y[n]=
\sum_{k=-\infty}^{\infty}
u[k]u[n-k].
\]

The terms are nonzero when

\[
k\ge0
\]

and

\[
k\le n.
\]

Hence,

\[
y[n]=
\sum_{k=0}^{n}1
=
n+1.
\]

Therefore,

\[
\boxed{
u[n]*u[n]=(n+1)u[n]
}
\]

%%%%%%%%%%%%%%%%%%%%%%%%%%%%%%%%%%%%%%%%%%%%%%%%%%%%%%%%%%%%%%%%%%%%%%%%%%%%%%%%

\section{Example 2}

Let

\[
x[n]=\delta[n]+\delta[n-1],
\]

\[
h[n]=\delta[n]+\delta[n-1].
\]

Using convolution,

\[
y[n]=
\delta[n]
+2\delta[n-1]
+\delta[n-2].
\]

%%%%%%%%%%%%%%%%%%%%%%%%%%%%%%%%%%%%%%%%%%%%%%%%%%%%%%%%%%%%%%%%%%%%%%%%%%%%%%%%

\section{Properties of Convolution}

%%%%%%%%%%%%%%%%%%%%%%%%%%%%%%%%%%%%%%%%%%%%%%%%%%%%%%%%%%%%%%%%%%%%%%%%%%%%%%%%

\subsection{Commutative}

\[
\boxed{
x[n]*h[n]=h[n]*x[n]
}
\]

%%%%%%%%%%%%%%%%%%%%%%%%%%%%%%%%%%%%%%%%%%%%%%%%%%%%%%%%%%%%%%%%%%%%%%%%%%%%%%%%

\subsection{Associative}

\[
\boxed{
(x_1*x_2)*x_3=x_1*(x_2*x_3)
}
\]

%%%%%%%%%%%%%%%%%%%%%%%%%%%%%%%%%%%%%%%%%%%%%%%%%%%%%%%%%%%%%%%%%%%%%%%%%%%%%%%%

\subsection{Distributive}

\[
\boxed{
x*(h_1+h_2)=x*h_1+x*h_2
}
\]

%%%%%%%%%%%%%%%%%%%%%%%%%%%%%%%%%%%%%%%%%%%%%%%%%%%%%%%%%%%%%%%%%%%%%%%%%%%%%%%%

\section{Difference Equation Representation}

Many discrete-time systems are described by difference equations.

%%%%%%%%%%%%%%%%%%%%%%%%%%%%%%%%%%%%%%%%%%%%%%%%%%%%%%%%%%%%%%%%%%%%%%%%%%%%%%%%

\subsection{First-Order System}

\[
\boxed{
y[n]-ay[n-1]=x[n]
}
\]

where \(a\) is a constant.

This is the discrete-time counterpart of a first-order differential equation.

%%%%%%%%%%%%%%%%%%%%%%%%%%%%%%%%%%%%%%%%%%%%%%%%%%%%%%%%%%%%%%%%%%%%%%%%%%%%%%%%

\section{Finding the Impulse Response}

Apply

\[
x[n]=\delta[n].
\]

Then

\[
h[n]-ah[n-1]=\delta[n].
\]

Solving recursively gives

\[
\boxed{
h[n]=a^n u[n]
}
\]

%%%%%%%%%%%%%%%%%%%%%%%%%%%%%%%%%%%%%%%%%%%%%%%%%%%%%%%%%%%%%%%%%%%%%%%%%%%%%%%%

\section{Memoryless Systems}

A system is memoryless if

\[
y[n]
\]

depends only on

\[
x[n].
\]

%%%%%%%%%%%%%%%%%%%%%%%%%%%%%%%%%%%%%%%%%%%%%%%%%%%%%%%%%%%%%%%%%%%%%%%%%%%%%%%%

\subsection{Example}

\[
y[n]=5x[n]
\]

is memoryless.

%%%%%%%%%%%%%%%%%%%%%%%%%%%%%%%%%%%%%%%%%%%%%%%%%%%%%%%%%%%%%%%%%%%%%%%%%%%%%%%%

\subsection{With Memory}

\[
y[n]=x[n]+x[n-1]
\]

has memory.

%%%%%%%%%%%%%%%%%%%%%%%%%%%%%%%%%%%%%%%%%%%%%%%%%%%%%%%%%%%%%%%%%%%%%%%%%%%%%%%%

\begin{tcolorbox}[title=Quick Recall,colback=blue!5]

\[
x[n]=x(nT_s)
\]

\[
\delta[n]=
\begin{cases}
1,& n=0\\
0,& n\neq0
\end{cases}
\]

\[
u[n]=
\begin{cases}
1,& n\ge0\\
0,& n<0
\end{cases}
\]

\[
y[n]=
\sum_{k=-\infty}^{\infty}
x[k]h[n-k]
\]

\[
u[n]*u[n]=(n+1)u[n]
\]

Causal:

\[
y[n]\text{ depends on }x[n],x[n-1],x[n-2],\ldots
\]

Stable:

\[
\sum |h[n]|<\infty
\]

\end{tcolorbox}
%%%%%%%%%%%%%%%%%%%%%%%%%%%%%%%%%%%%%%%%%%%%%%%%%%%%%%%%%%%%%%%%%%%%%%%%%%%%%%%%
%%%%%%%%%%%%%%%%%%%%%%%%%%%%%%%%%%%%%%%%%%%%%%%%%%%%%%%%%%%%%%%%%%%%%%%%%%%%%%%%
\section{Discrete-Time Fourier Transform (DTFT)}

For an aperiodic discrete-time sequence \(x[n]\), the frequency-domain
representation is the \textbf{Discrete-Time Fourier Transform (DTFT)}.

%%%%%%%%%%%%%%%%%%%%%%%%%%%%%%%%%%%%%%%%%%%%%%%%%%%%%%%%%%%%%%%%%%%%%%%%%%%%%%%%

\subsection{Definition}

\[
\boxed{
X(e^{j\omega})=
\sum_{n=-\infty}^{\infty}
x[n]e^{-j\omega n}
}
\]

where

\[
-\pi\le\omega\le\pi.
\]

%%%%%%%%%%%%%%%%%%%%%%%%%%%%%%%%%%%%%%%%%%%%%%%%%%%%%%%%%%%%%%%%%%%%%%%%%%%%%%%%

\subsection{Inverse DTFT}

\[
\boxed{
x[n]=
\frac{1}{2\pi}
\int_{-\pi}^{\pi}
X(e^{j\omega})e^{j\omega n}\,d\omega
}
\]

The DTFT is periodic in frequency with period \(2\pi\):

\[
\boxed{
X(e^{j(\omega+2\pi)})=X(e^{j\omega})
}
\]

%%%%%%%%%%%%%%%%%%%%%%%%%%%%%%%%%%%%%%%%%%%%%%%%%%%%%%%%%%%%%%%%%%%%%%%%%%%%%%%%

\section{DTFT of Basic Sequences}

%%%%%%%%%%%%%%%%%%%%%%%%%%%%%%%%%%%%%%%%%%%%%%%%%%%%%%%%%%%%%%%%%%%%%%%%%%%%%%%%

\subsection{Unit Sample}

\[
\boxed{
\delta[n]\xleftrightarrow{\mathcal{F}}1
}
\]

%%%%%%%%%%%%%%%%%%%%%%%%%%%%%%%%%%%%%%%%%%%%%%%%%%%%%%%%%%%%%%%%%%%%%%%%%%%%%%%%

\subsection{Shifted Impulse}

\[
\boxed{
\delta[n-n_0]
\xleftrightarrow{\mathcal{F}}
e^{-j\omega n_0}
}
\]

%%%%%%%%%%%%%%%%%%%%%%%%%%%%%%%%%%%%%%%%%%%%%%%%%%%%%%%%%%%%%%%%%%%%%%%%%%%%%%%%

\subsection{Unit Step}

For \(|a|<1\),

\[
x[n]=a^n u[n]
\]

has DTFT

\[
\boxed{
X(e^{j\omega})=
\frac{1}{1-ae^{-j\omega}}
}
\]

This is one of the most important transform pairs.

%%%%%%%%%%%%%%%%%%%%%%%%%%%%%%%%%%%%%%%%%%%%%%%%%%%%%%%%%%%%%%%%%%%%%%%%%%%%%%%%

\section{Frequency Response of an LTI System}

For an LTI system with impulse response \(h[n]\),

\[
\boxed{
H(e^{j\omega})=
\sum_{n=-\infty}^{\infty}
h[n]e^{-j\omega n}
}
\]

The output spectrum is

\[
\boxed{
Y(e^{j\omega})=
X(e^{j\omega})H(e^{j\omega})
}
\]

%%%%%%%%%%%%%%%%%%%%%%%%%%%%%%%%%%%%%%%%%%%%%%%%%%%%%%%%%%%%%%%%%%%%%%%%%%%%%%%%

\section{Example: Moving Average Filter}

Consider

\[
h[n]=
\begin{cases}
1, & n=0,1\\
0, & \text{otherwise}
\end{cases}
\]

Then

\[
H(e^{j\omega})=
1+e^{-j\omega}
\]

Using Euler’s identity,

\[
H(e^{j\omega})=
2e^{-j\omega/2}\cos\frac{\omega}{2}
\]

Hence,

\[
\boxed{
|H(e^{j\omega})|=2\left|\cos\frac{\omega}{2}\right|
}
\]

and

\[
\boxed{
\angle H(e^{j\omega})=-\frac{\omega}{2}
}
\]

This is a simple low-pass filter.

%%%%%%%%%%%%%%%%%%%%%%%%%%%%%%%%%%%%%%%%%%%%%%%%%%%%%%%%%%%%%%%%%%%%%%%%%%%%%%%%

\section{Difference Equation}

A general linear constant-coefficient difference equation is

\[
\boxed{
\sum_{k=0}^{N}a_k y[n-k]
=
\sum_{k=0}^{M}b_k x[n-k]
}
\]

Taking the DTFT,

\[
\left(
\sum_{k=0}^{N}a_k e^{-j\omega k}
\right)
Y(e^{j\omega})
=
\left(
\sum_{k=0}^{M}b_k e^{-j\omega k}
\right)
X(e^{j\omega})
\]

Therefore,

\[
\boxed{
H(e^{j\omega})=
\frac{\sum_{k=0}^{M}b_k e^{-j\omega k}}
{\sum_{k=0}^{N}a_k e^{-j\omega k}}
}
\]

%%%%%%%%%%%%%%%%%%%%%%%%%%%%%%%%%%%%%%%%%%%%%%%%%%%%%%%%%%%%%%%%%%%%%%%%%%%%%%%%

\section{z-Transform}

The DTFT exists only for some sequences. A more general transform is the
\textbf{z-transform}.

%%%%%%%%%%%%%%%%%%%%%%%%%%%%%%%%%%%%%%%%%%%%%%%%%%%%%%%%%%%%%%%%%%%%%%%%%%%%%%%%

\subsection{Definition}

\[
\boxed{
X(z)=
\sum_{n=-\infty}^{\infty}
x[n]z^{-n}
}
\]

where

\[
z=re^{j\omega}.
\]

%%%%%%%%%%%%%%%%%%%%%%%%%%%%%%%%%%%%%%%%%%%%%%%%%%%%%%%%%%%%%%%%%%%%%%%%%%%%%%%%

\subsection{Region of Convergence (ROC)}

The set of values of \(z\) for which the series converges is called the ROC.

The ROC is essential for determining causality and stability.

%%%%%%%%%%%%%%%%%%%%%%%%%%%%%%%%%%%%%%%%%%%%%%%%%%%%%%%%%%%%%%%%%%%%%%%%%%%%%%%%

\section{z-Transform of \(a^n u[n]\)}

\[
X(z)=
\sum_{n=0}^{\infty}a^n z^{-n}
\]

This is a geometric series:

\[
X(z)=
\frac{1}{1-az^{-1}}
\]

with ROC

\[
\boxed{
|z|>|a|
}
\]

%%%%%%%%%%%%%%%%%%%%%%%%%%%%%%%%%%%%%%%%%%%%%%%%%%%%%%%%%%%%%%%%%%%%%%%%%%%%%%%%

\section{Poles and Zeros}

%%%%%%%%%%%%%%%%%%%%%%%%%%%%%%%%%%%%%%%%%%%%%%%%%%%%%%%%%%%%%%%%%%%%%%%%%%%%%%%%

\subsection{Zero}

A value of \(z\) for which

\[
X(z)=0.
\]

%%%%%%%%%%%%%%%%%%%%%%%%%%%%%%%%%%%%%%%%%%%%%%%%%%%%%%%%%%%%%%%%%%%%%%%%%%%%%%%%

\subsection{Pole}

A value of \(z\) for which

\[
X(z)\to\infty.
\]

%%%%%%%%%%%%%%%%%%%%%%%%%%%%%%%%%%%%%%%%%%%%%%%%%%%%%%%%%%%%%%%%%%%%%%%%%%%%%%%%

\subsection{Example}

\[
X(z)=
\frac{1-z^{-1}}
{1-0.5z^{-1}}
\]

Zero:

\[
z=1
\]

Pole:

\[
z=0.5
\]

%%%%%%%%%%%%%%%%%%%%%%%%%%%%%%%%%%%%%%%%%%%%%%%%%%%%%%%%%%%%%%%%%%%%%%%%%%%%%%%%

\section{Pole-Zero Plot}

The pole-zero plot is drawn in the complex \(z\)-plane.

\begin{center}
\begin{tikzpicture}[scale=2]
\draw[->] (-1.5,0) -- (1.5,0) node[right] {Re};
\draw[->] (0,-1.2) -- (0,1.2) node[above] {Im};

\draw (0,0) circle (1);

\draw[thick] (1,0) node[circle,draw,inner sep=1.5pt] {};
\draw[thick] (0.5,0) node[cross out,draw,inner sep=2pt] {};

\node[below] at (1,0) {$z=1$};
\node[below] at (0.5,0) {$z=0.5$};
\end{tikzpicture}
\end{center}

Convention:

\begin{itemize}
\item \(\circ\): zero,
\item \(\times\): pole.
\end{itemize}

%%%%%%%%%%%%%%%%%%%%%%%%%%%%%%%%%%%%%%%%%%%%%%%%%%%%%%%%%%%%%%%%%%%%%%%%%%%%%%%%

\section{Causality from ROC}

A rational system is causal if the ROC extends outward from the outermost pole.

Example:

\[
X(z)=
\frac{1}{1-0.5z^{-1}}
\]

ROC:

\[
|z|>0.5
\]

Hence,

\[
\boxed{
\text{The sequence is causal.}
}
\]

%%%%%%%%%%%%%%%%%%%%%%%%%%%%%%%%%%%%%%%%%%%%%%%%%%%%%%%%%%%%%%%%%%%%%%%%%%%%%%%%

\section{Stability from ROC}

A system is BIBO stable if the ROC includes the unit circle

\[
\boxed{
|z|=1
}
\]

For the previous example,

\[
|z|>0.5
\]

includes the unit circle, so

\[
\boxed{
\text{The system is stable.}
}
\]

%%%%%%%%%%%%%%%%%%%%%%%%%%%%%%%%%%%%%%%%%%%%%%%%%%%%%%%%%%%%%%%%%%%%%%%%%%%%%%%%

\section{Frequency Response from z-Transform}

The DTFT is obtained by evaluating the z-transform on the unit circle:

\[
\boxed{
H(e^{j\omega})=H(z)\big|_{z=e^{j\omega}}
}
\]

Thus, the unit circle is the frequency axis of discrete-time systems.

%%%%%%%%%%%%%%%%%%%%%%%%%%%%%%%%%%%%%%%%%%%%%%%%%%%%%%%%%%%%%%%%%%%%%%%%%%%%%%%%

\section{First-Order IIR System}

Consider

\[
y[n]-0.5y[n-1]=x[n].
\]

Taking the z-transform,

\[
Y(z)-0.5z^{-1}Y(z)=X(z).
\]

Hence,

\[
H(z)=
\frac{Y(z)}{X(z)}
=
\frac{1}{1-0.5z^{-1}}.
\]

%%%%%%%%%%%%%%%%%%%%%%%%%%%%%%%%%%%%%%%%%%%%%%%%%%%%%%%%%%%%%%%%%%%%%%%%%%%%%%%%

\subsection{Pole}

\[
z=0.5
\]

%%%%%%%%%%%%%%%%%%%%%%%%%%%%%%%%%%%%%%%%%%%%%%%%%%%%%%%%%%%%%%%%%%%%%%%%%%%%%%%%

\subsection{Frequency Response}

\[
\boxed{
H(e^{j\omega})=
\frac{1}{1-0.5e^{-j\omega}}
}
\]

%%%%%%%%%%%%%%%%%%%%%%%%%%%%%%%%%%%%%%%%%%%%%%%%%%%%%%%%%%%%%%%%%%%%%%%%%%%%%%%%

\section{Magnitude Response}

\[
|H(e^{j\omega})|=
\frac{1}
{\sqrt{1+0.25-\cos\omega}}
\]

%%%%%%%%%%%%%%%%%%%%%%%%%%%%%%%%%%%%%%%%%%%%%%%%%%%%%%%%%%%%%%%%%%%%%%%%%%%%%%%%

\section{Phase Response}

\[
\angle H(e^{j\omega})=
-\tan^{-1}
\left(
\frac{0.5\sin\omega}
{1-0.5\cos\omega}
\right)
\]

%%%%%%%%%%%%%%%%%%%%%%%%%%%%%%%%%%%%%%%%%%%%%%%%%%%%%%%%%%%%%%%%%%%%%%%%%%%%%%%%

\section{Phase Delay}

\[
\boxed{
\tau_p(\omega)=
-\frac{\angle H(e^{j\omega})}{\omega}
}
\]

Phase delay measures the delay experienced by a sinusoid of frequency \(\omega\).

%%%%%%%%%%%%%%%%%%%%%%%%%%%%%%%%%%%%%%%%%%%%%%%%%%%%%%%%%%%%%%%%%%%%%%%%%%%%%%%%

\section{Group Delay}

\[
\boxed{
\tau_g(\omega)=
-\frac{d}{d\omega}\angle H(e^{j\omega})
}
\]

Group delay measures the delay experienced by the envelope of a narrowband signal.

%%%%%%%%%%%%%%%%%%%%%%%%%%%%%%%%%%%%%%%%%%%%%%%%%%%%%%%%%%%%%%%%%%%%%%%%%%%%%%%%

\section{Linear Phase}

A system has linear phase if

\[
\angle H(e^{j\omega})=-\omega n_0.
\]

Then,

\[
\tau_p(\omega)=\tau_g(\omega)=n_0.
\]

Linear-phase systems do not distort waveform shape.

%%%%%%%%%%%%%%%%%%%%%%%%%%%%%%%%%%%%%%%%%%%%%%%%%%%%%%%%%%%%%%%%%%%%%%%%%%%%%%%%

\section{FIR vs IIR Systems}

\begin{table}[H]
\centering
\caption{Comparison of FIR and IIR systems}
\begin{tabular}{|p{5cm}|p{3cm}|p{3cm}|}
\hline
\textbf{Feature} & \textbf{FIR} & \textbf{IIR}\\
\hline
Impulse response & Finite & Infinite\\
\hline
Poles & At origin only & General poles\\
\hline
Always stable? & Yes & Not always\\
\hline
Linear phase possible? & Yes & Generally no\\
\hline
Computation & More multipliers & Fewer multipliers\\
\hline
\end{tabular}
\end{table}

%%%%%%%%%%%%%%%%%%%%%%%%%%%%%%%%%%%%%%%%%%%%%%%%%%%%%%%%%%%%%%%%%%%%%%%%%%%%%%%%

\section{Solved Example 7.1}

Find the DTFT of

\[
x[n]=\left(\frac12\right)^n u[n].
\]

%%%%%%%%%%%%%%%%%%%%%%%%%%%%%%%%%%%%%%%%%%%%%%%%%%%%%%%%%%%%%%%%%%%%%%%%%%%%%%%%

\subsection{Solution}

Using the standard pair,

\[
\boxed{
X(e^{j\omega})=
\frac{1}{1-\frac12 e^{-j\omega}}
}
\]

%%%%%%%%%%%%%%%%%%%%%%%%%%%%%%%%%%%%%%%%%%%%%%%%%%%%%%%%%%%%%%%%%%%%%%%%%%%%%%%%

\section{Solved Example 7.2}

Determine whether the system

\[
y[n]=x[n]+x[n-1]
\]

is causal and stable.

%%%%%%%%%%%%%%%%%%%%%%%%%%%%%%%%%%%%%%%%%%%%%%%%%%%%%%%%%%%%%%%%%%%%%%%%%%%%%%%%

\subsection{Impulse Response}

\[
h[n]=\delta[n]+\delta[n-1]
\]

%%%%%%%%%%%%%%%%%%%%%%%%%%%%%%%%%%%%%%%%%%%%%%%%%%%%%%%%%%%%%%%%%%%%%%%%%%%%%%%%

\subsection{Causality}

Both impulses occur at \(n=0\) and \(n=1\), so the system is causal.

%%%%%%%%%%%%%%%%%%%%%%%%%%%%%%%%%%%%%%%%%%%%%%%%%%%%%%%%%%%%%%%%%%%%%%%%%%%%%%%%

\subsection{Stability}

\[
\sum |h[n]|=1+1=2<\infty
\]

Hence,

\[
\boxed{
\text{The system is causal and stable.}
}
\]

%%%%%%%%%%%%%%%%%%%%%%%%%%%%%%%%%%%%%%%%%%%%%%%%%%%%%%%%%%%%%%%%%%%%%%%%%%%%%%%%

\section{Solved Example 7.3}

Find the output of the system with

\[
h[n]=u[n]
\]

for input

\[
x[n]=u[n].
\]

%%%%%%%%%%%%%%%%%%%%%%%%%%%%%%%%%%%%%%%%%%%%%%%%%%%%%%%%%%%%%%%%%%%%%%%%%%%%%%%%

\subsection{Solution}

\[
y[n]=u[n]*u[n]
\]

\[
y[n]=(n+1)u[n]
\]

%%%%%%%%%%%%%%%%%%%%%%%%%%%%%%%%%%%%%%%%%%%%%%%%%%%%%%%%%%%%%%%%%%%%%%%%%%%%%%%%

\section{Solved Example 7.4}

For

\[
H(z)=
\frac{1}{1-1.2z^{-1}},
\]

determine stability.

%%%%%%%%%%%%%%%%%%%%%%%%%%%%%%%%%%%%%%%%%%%%%%%%%%%%%%%%%%%%%%%%%%%%%%%%%%%%%%%%

\subsection{Pole}

\[
z=1.2
\]

For a causal system, ROC is

\[
|z|>1.2
\]

The unit circle \(|z|=1\) is not included.

Therefore,

\[
\boxed{
\text{The system is unstable.}
}
\]

%%%%%%%%%%%%%%%%%%%%%%%%%%%%%%%%%%%%%%%%%%%%%%%%%%%%%%%%%%%%%%%%%%%%%%%%%%%%%%%%

\section{Quick Test}

\begin{enumerate}
\item \(x[n]=\cos(\pi n/3)\) has period \(\boxed{6}\).
\item \(u[n]*u[n]=\boxed{(n+1)u[n]}\).
\item Stable LTI system: \(\boxed{\sum |h[n]|<\infty}\).
\item DTFT periodicity: \(\boxed{2\pi}\).
\item Frequency response: \(\boxed{H(e^{j\omega})=H(z)|_{z=e^{j\omega}}}\).
\end{enumerate}

%%%%%%%%%%%%%%%%%%%%%%%%%%%%%%%%%%%%%%%%%%%%%%%%%%%%%%%%%%%%%%%%%%%%%%%%%%%%%%%%

\section{Chapter Summary}

In this chapter, we introduced discrete-time signals and systems and developed
the DTFT for frequency-domain analysis.

We studied

\begin{itemize}
\item sequence operations,
\item periodicity,
\item system properties,
\item convolution,
\item DTFT,
\item z-transform,
\item poles and zeros,
\item causality and stability,
\item and frequency response.
\end{itemize}

We also introduced phase delay and group delay, which are important in digital
filter design and communication systems.

%%%%%%%%%%%%%%%%%%%%%%%%%%%%%%%%%%%%%%%%%%%%%%%%%%%%%%%%%%%%%%%%%%%%%%%%%%%%%%%%

\begin{tcolorbox}[title=One-Minute Revision,colback=blue!5]

\[
X(e^{j\omega})=
\sum_{n=-\infty}^{\infty}x[n]e^{-j\omega n}
\]

\[
x[n]=
\frac{1}{2\pi}
\int_{-\pi}^{\pi}
X(e^{j\omega})e^{j\omega n}d\omega
\]

\[
X(e^{j\omega})\text{ is periodic with }2\pi
\]

\[
y[n]=
\sum_{k=-\infty}^{\infty}
x[k]h[n-k]
\]

\[
X(z)=
\sum_{n=-\infty}^{\infty}x[n]z^{-n}
\]

\[
H(e^{j\omega})=H(z)|_{z=e^{j\omega}}
\]

Stable:

\[
\sum |h[n]|<\infty
\]

Causal ROC:

\[
|z|>\text{outermost pole}
\]

Stable ROC includes:

\[
|z|=1
\]

\[
\tau_g(\omega)=
-\frac{d}{d\omega}\angle H(e^{j\omega})
\]

\end{tcolorbox}
%%%%%%%%%%%%%%%%%%%%%%%%%%%%%%%%%%%%%%%%%%%%%%%%%%%%%%%%%%%%%%%%%%%%%%%%%%%%%%%%